\section{The Littlewood--Richardson rule}

\subsection{Tuple arithmetic}

\begin{definition}
\label{def.sf.tuple-addition}
\lean{AlgebraicCombinatorics.zeroTuple, AlgebraicCombinatorics.addTuple, AlgebraicCombinatorics.subTuple}
\leantarget
\leanok
\textbf{(a)} We let $\mathbf{0}$ denote the
$N$-tuple $\left(  0,0,\ldots,0\right)  \in\mathbb{N}^{N}$.

\medskip

\textbf{(b)} Let $\alpha=\left(  \alpha_{1},\alpha_{2},\ldots,\alpha
_{N}\right)  $ and $\beta=\left(  \beta_{1},\beta_{2},\ldots,\beta_{N}\right)
$ be two $N$-tuples in $\mathbb{N}^{N}$. Then, we set%
\begin{align*}
\alpha+\beta &  :=\left(  \alpha_{1}+\beta_{1},\alpha_{2}+\beta_{2}%
,\ldots,\alpha_{N}+\beta_{N}\right)  \ \ \ \ \ \ \ \ \ \ \text{and}\\
\alpha-\beta &  :=\left(  \alpha_{1}-\beta_{1},\alpha_{2}-\beta_{2}%
,\ldots,\alpha_{N}-\beta_{N}\right)  .
\end{align*}
Note that $\alpha+\beta\in\mathbb{N}^{N}$, whereas $\alpha-\beta\in
\mathbb{Z}^{N}$.
\end{definition}

\subsection{Content of a tableau}

\begin{definition}
\label{def.sf.content}
\lean{AlgebraicCombinatorics.contentTableau}
\leantarget
\leanok
Let $\lambda$ and $\mu$ be two $N$-partitions. Let $T$
be a tableau of shape $\lambda/\mu$. We define the \emph{content} of $T$ to be
the $N$-tuple $\left(  a_{1},a_{2},\ldots,a_{N}\right)  $, where%
\[
a_{i}:=\left(  \text{\# of }i\text{'s in }T\right)  =\left(  \text{\# of boxes
}c\text{ of }T\text{ such that }T\left(  c\right)  =i\right)  .
\]
We denote this $N$-tuple by $\operatorname*{cont}T$.
\end{definition}

\subsection{Column restriction of tableaux}

\begin{definition}
\label{def.sf.col-tab}
\lean{AlgebraicCombinatorics.colGeq}
\leantarget
\leanok
Let $\lambda$ and $\mu$ be two $N$-partitions. Let $T$
be a tableau of shape $\lambda/\mu$. Let $j$ be a positive integer. Then,
$\operatorname{col}_{\geq j}T$ means the restriction of $T$ to columns
$j,j+1,j+2,\ldots$ (that is, the result of removing the first $j-1$ columns
from $T$). Formally speaking, this means the restriction of the map $T$ to the
set $\left\{  \left(  u,v\right)  \in Y\left(  \lambda/\mu\right)
\ \mid\ v\geq j\right\}  $.
\end{definition}

\subsection{Yamanouchi tableaux}

\begin{definition}
\label{def.sf.yamanouchi}
\lean{AlgebraicCombinatorics.IsYamanouchi}
\leantarget
\leanok
Let $\lambda,\mu,\nu$ be three $N$-partitions. A
semistandard tableau $T$ of shape $\lambda/\mu$ is said to be $\nu
$-\emph{Yamanouchi} (this is an adjective) if for each positive integer $j$,
the $N$-tuple $\nu+\operatorname*{cont}\left(  \operatorname{col}_{\geq
j}T\right)  \in\mathbb{N}^{N}$ is an $N$-partition (i.e., weakly decreasing).
\end{definition}

\subsection{Content and column restriction helpers}

\begin{lemma}[Content of column-restricted tableau]
\label{lem.sf.content-col-geq-one}
\lean{AlgebraicCombinatorics.contentColGeq_one_eq_contentTableau}
\leanhelper
\leanok
Let $T$ be a tableau of shape $\lambda/\mu$. Then
$\operatorname{cont}(\operatorname{col}_{\geq 1} T) = \operatorname{cont}(T)$.
In other words, since every cell in $Y(\lambda/\mu)$ has column $\geq 1$,
the column restriction to $j = 1$ recovers the full content.
\end{lemma}

\begin{proof}
\leanok
Every cell $(i, v) \in Y(\lambda/\mu)$ satisfies $v > \mu_i \geq 0$,
hence $v \geq 1$, so $\operatorname{col}_{\geq 1} T$ has the same domain as~$T$.
\end{proof}

\begin{lemma}[Content monotonicity under column restriction]
\label{lem.sf.content-col-geq-mono}
\lean{AlgebraicCombinatorics.contentColGeq_mono}
\leanhelper
\leanok
For any tableau $T$ of shape $\lambda/\mu$, row index $i$, and positive
integer $j$:
\[
  \operatorname{cont}(\operatorname{col}_{\geq j+1} T)_i
  \;\leq\;
  \operatorname{cont}(\operatorname{col}_{\geq j} T)_i.
\]
That is, restricting to higher columns can only decrease the count
of any entry.
\end{lemma}

\begin{proof}
\leanok
The cells counted by $\operatorname{cont}(\operatorname{col}_{\geq j+1} T)_i$
form a subset of those counted by $\operatorname{cont}(\operatorname{col}_{\geq j} T)_i$.
\end{proof}

\begin{lemma}[Column restriction vanishes beyond the shape]
\label{lem.sf.content-col-geq-zero}
\lean{AlgebraicCombinatorics.contentColGeq_eq_zero_of_large}
\leanhelper
\leanok
If $j > \lambda_i$ for every row index $i$, then
$\operatorname{cont}(\operatorname{col}_{\geq j} T) = \mathbf{0}$.
\end{lemma}

\begin{proof}
\leanok
When $j$ exceeds all column lengths, the domain
$\{(u,v) \in Y(\lambda/\mu) \mid v \geq j\}$ is empty.
\end{proof}

\begin{lemma}[Content equality from pointwise characterization]
\label{lem.sf.content-eq-of-iff}
\lean{AlgebraicCombinatorics.contentTableau_eq_of_iff}
\leanhelper
\leanok
If two tableaux $T, T'$ of the same shape satisfy
$T'(c) = i \iff T(c) = i$ for all cells $c$, then
$\operatorname{cont}(T')_i = \operatorname{cont}(T)_i$.
\end{lemma}

\begin{proof}
\leanok
The sets $\{c \mid T'(c) = i\}$ and $\{c \mid T(c) = i\}$ are equal
by the hypothesis (each cell belongs to one set if and only if it belongs to the other),
so they have the same cardinality.
\end{proof}

\subsection{Properties of alternants}

\begin{lemma}
\label{lem.sf.alternant-0-lr}
\lean{AlgebraicCombinatorics.alternant_eq_zero_of_repeated, AlgebraicCombinatorics.alternant_swap}
\leantarget
\leanok
Let $\alpha\in\mathbb{N}^{N}$.

\medskip

\textbf{(a)} If the $N$-tuple $\alpha$ has two equal entries, then $a_{\alpha
}=0$.

\medskip

\textbf{(b)} Let $\beta\in\mathbb{N}^{N}$ be an $N$-tuple obtained from
$\alpha$ by swapping two entries. Then, $a_{\beta}=-a_{\alpha}$.
\end{lemma}

\begin{proof}
\leanok
Both parts follow from the definition of alternants as signed sums over
permutations.

\textbf{(a)} If $\alpha_i = \alpha_j$ with $i \neq j$, then the map
$\sigma \mapsto (i\;j) \circ \sigma$ is a sign-reversing involution on $S_N$
that preserves $x^{\alpha \circ \sigma}$ (since $\alpha$ is invariant under
swapping $i$ and $j$). The terms cancel in pairs, giving $a_{\alpha} = 0$.

\textbf{(b)} If $\beta = \alpha \circ (i\;j)$, then reindexing the sum
defining $a_{\beta}$ via $\tau \mapsto (i\;j) \circ \tau$ introduces a factor
of $\operatorname{sign}((i\;j)) = -1$, giving $a_{\beta} = -a_{\alpha}$.
\end{proof}

\subsection{Regular elements and cancellation}

\begin{definition}
\label{def.cring.reg}
\lean{AlgebraicCombinatorics.IsRegularElement}
\leantarget
\leanok
Let $L$ be a commutative ring. Let $a\in L$. The element
$a$ of $L$ is said to be \emph{regular} if and only if every $x\in L$
satisfying $ax=0$ satisfies $x=0$.
\end{definition}

\begin{lemma}
\label{lem.cring.reg.cancel}
\lean{AlgebraicCombinatorics.IsRegularElement.cancel}
\leantarget
\leanok
Let $L$ be a commutative ring. Let $a,u,v\in L$ be
such that $a$ is regular. Assume that $au=av$. Then, $u=v$.
\end{lemma}

\begin{proof}
\leanok
We have $a\left(  u-v\right)
=au-av=0$ (since $au=av$). However, $a$ is regular; in other words, every
$x\in L$ satisfying $ax=0$ satisfies $x=0$ (by the definition of
``regular''). Applying this to $x=u-v$, we
obtain $u-v=0$ (since $a\left(  u-v\right)  =0$). Thus, $u=v$.
\end{proof}

\begin{lemma}
\label{lem.sf.arho-reg}
\lean{AlgebraicCombinatorics.alternant_rho_isRegular}
\leantarget
\leanok
The element $a_{\rho}$ of the polynomial ring
$\mathcal{P}$ is regular.
\end{lemma}

\begin{proof}
\leanok
Each of the
indeterminates $x_{1},x_{2},\ldots,x_{N}$ is regular (as an element of
$\mathcal{P}$). Indeed, multiplying a polynomial $f$ by an indeterminate
$x_{i}$ merely shifts the coefficients of $f$ to different monomials; thus, if
$x_{i}f=0$, then $f=0$. Next, the polynomial $x_{i}-x_{j}%
\in\mathcal{P}$ is regular whenever $1\leq i<j\leq N$. Indeed, this polynomial
$x_{i}-x_{j}$ is the image of the indeterminate $x_{i}$ under a
$K$-algebra automorphism of $\mathcal{P}$ (the automorphism that
sends $x_{i}$ to $x_{i}-x_{j}$ while leaving all other indeterminates
unchanged), and therefore is regular because $x_{i}$ is
regular (and because any ring automorphism sends regular elements to regular
elements). Any finite product of regular elements is again regular. Thus,
the element $\prod_{1\leq i<j\leq N}\left(  x_{i}-x_{j}\right)  \in
\mathcal{P}$ is regular. Since
$a_{\rho} = \prod_{1\leq i<j\leq N}\left(  x_{i}-x_{j}\right)$
(the Vandermonde determinant), the element
$a_{\rho}\in\mathcal{P}$ is regular.

\emph{Lean note:} The Lean formalization takes a shortcut by assuming $R$ is
an integral domain, which makes the polynomial ring an integral domain;
regularity then follows from the Vandermonde polynomial being nonzero.
\end{proof}

\subsection{Bender--Knuth involution}

\begin{definition}[Forced and free cells]
\label{def.sf.bk-forced-free}
\lean{AlgebraicCombinatorics.isForcedK, AlgebraicCombinatorics.isForcedKSucc, AlgebraicCombinatorics.isUnmatchedFreeK, AlgebraicCombinatorics.isUnmatchedFreeKSucc}
\leanhelper
\leanok
Let $T$ be a semistandard tableau and let $k < N-1$.
A cell $c = (i,j)$ with $T(c) = k$ is \emph{forced} if there exists a cell
$(i+1,j)$ directly below with $T(i+1,j) = k+1$. Similarly for $T(c) = k+1$ with a cell directly above
having value~$k$. Free cells are those that are not forced.

Among the free cells, the \emph{parenthesis-matching algorithm} classifies
each as \emph{matched} or \emph{unmatched}: reading free entries left-to-right
in each row, each free $(k{+}1)$ matches with the nearest unmatched free $k$
to its left. Unmatched entries are those remaining after this process.
\end{definition}

\begin{definition}[Bender--Knuth involution $\beta_k$]
\label{def.sf.bk-involution}
\lean{AlgebraicCombinatorics.benderKnuth}
\leanhelper
\leanok
Let $T$ be a semistandard tableau of shape $\lambda/\mu$ and let
$k < N-1$.  The \emph{Bender--Knuth involution} $\beta_k(T)$ is obtained
from $T$ by swapping only the \emph{unmatched} free $k$'s and
unmatched free $(k{+}1)$'s:
\[
  \beta_k(T)(c) = \begin{cases}
    k+1 & \text{if } c \text{ is an unmatched free } k, \\
    k   & \text{if } c \text{ is an unmatched free } k{+}1, \\
    T(c) & \text{otherwise.}
  \end{cases}
\]
Forced cells and matched free cells are left unchanged.
\end{definition}

\begin{lemma}[Bender--Knuth preserves semistandardness]
\label{lem.sf.bk-semistandard}
\lean{AlgebraicCombinatorics.benderKnuth_semistandard}
\leanhelper
\leanok
If $T$ is semistandard and $\lambda$, $\mu$ are $N$-partitions, then
$\beta_k(T)$ is semistandard.
\end{lemma}

\begin{proof}
\leanok
Row-weak ordering is preserved because in each row, unmatched $(k{+}1)$'s
(which become $k$) lie to the left of unmatched $k$'s (which become $k{+}1$),
by the structure of the parenthesis matching (Lemma~\ref{lem.sf.bk-row-weak}).
Column-strict ordering is preserved because the dangerous case
(an unmatched $k$ above an unmatched $k{+}1$ in the same column) is ruled out
by the partition property, which forces adjacent $k, k{+}1$ pairs to be forced
(Lemma~\ref{lem.sf.bk-column-strict}).
\end{proof}

\begin{lemma}[Bender--Knuth row-weak preservation]
\label{lem.sf.bk-row-weak}
\lean{AlgebraicCombinatorics.benderKnuth_row_weak}
\leanhelper
\leanok
For cells $(i, j_1)$ and $(i, j_2)$ in the same row with $j_1 < j_2$,
$\beta_k(T)(i, j_1) \leq \beta_k(T)(i, j_2)$.
\end{lemma}

\begin{proof}
\leanok
The proof proceeds by case analysis on the status of $c_1$ and $c_2$
(forced, matched free, unmatched free, or having value other than $k$ or $k{+}1$).
The key insight is that unmatched $(k{+}1)$'s are the \emph{leftmost} free
$(k{+}1)$'s and unmatched $k$'s are the \emph{rightmost} free $k$'s in each row.
\end{proof}

\begin{lemma}[Bender--Knuth column-strict preservation]
\label{lem.sf.bk-column-strict}
\lean{AlgebraicCombinatorics.benderKnuth_column_strict}
\leanhelper
\leanok
For cells $(i_1, j)$ and $(i_2, j)$ in the same column with $i_1 < i_2$,
$\beta_k(T)(i_1, j) < \beta_k(T)(i_2, j)$.
\end{lemma}

\begin{proof}
\leanok
The proof proceeds by case analysis on whether $c_1$ and $c_2$ are
unmatched free $k$, unmatched free $(k{+}1)$, or unchanged.
The dangerous case is when $c_1$ (above) is an unmatched free $k$
(becoming $k{+}1$) and $c_2$ (below) is an unmatched free $(k{+}1)$
(becoming $k$), which would violate strict ordering. But this cannot happen:
if $T(c_1) = k$ and $T(c_2) = k{+}1$ in the same column, the partition
property forces them to be in adjacent rows, making them a forced pair
(not free), so they are unchanged by $\beta_k$.
All other case combinations preserve strict ordering.
\end{proof}

\begin{lemma}[Bender--Knuth is an involution]
\label{lem.sf.bk-involutive}
\lean{AlgebraicCombinatorics.benderKnuth_involutive}
\leanhelper
\leanok
For $N$-partitions $\lambda$, $\mu$ and a semistandard tableau $T$ of shape
$\lambda/\mu$, the tableau $T' = \beta_k(T)$ is semistandard, and $\beta_k(T') = T$.
\end{lemma}

\begin{proof}
\leanok
The matching structure is symmetric under BK: an unmatched free $k$ becomes
an unmatched free $(k{+}1)$ in $T'$ and vice versa; matched and forced cells
are unchanged. Applying BK twice therefore returns every cell to its original value.
\end{proof}

\begin{lemma}[Forced cells are preserved under BK]
\label{lem.sf.bk-forced-preserved}
\lean{AlgebraicCombinatorics.forced_k_preserved', AlgebraicCombinatorics.forced_kSucc_preserved'}
\leanhelper
\leanok
If $c$ is a forced $k$ cell in $T$, then $c$ is also a forced $k$ cell in
$\beta_k(T)$. Similarly for forced $(k{+}1)$ cells.
\end{lemma}

\begin{proof}
\leanok
A forced $k$ cell has a forced $(k{+}1)$ partner directly below in the same
column. Since forced cells are never unmatched, both cells retain their values
under $\beta_k$, preserving the forced structure.
\end{proof}

\begin{lemma}[Unmatched cells swap status under BK]
\label{lem.sf.bk-unmatched-becomes-unmatched}
\lean{AlgebraicCombinatorics.unmatched_k_becomes_unmatched_kSucc', AlgebraicCombinatorics.unmatched_kSucc_becomes_unmatched_k'}
\leanhelper
\leanok
An unmatched free $k$ in $T$ becomes an unmatched free $(k{+}1)$ in
$\beta_k(T)$, and vice versa.
\end{lemma}

\begin{proof}
\leanok
After swapping the value from $k$ to $k{+}1$, the cell is free (not forced)
because there is no cell directly above with value $k$. The counting argument
shows that the cumulative counts shift appropriately to make the cell unmatched
in its new role.
\end{proof}

\begin{lemma}[Matched cells stay matched under BK]
\label{lem.sf.bk-matched-stays-matched}
\lean{AlgebraicCombinatorics.matched_k_stays_matched', AlgebraicCombinatorics.matched_kSucc_stays_matched'}
\leanhelper
\leanok
A matched free $k$ in $T$ remains a matched free $k$ in $\beta_k(T)$,
and similarly for matched free $(k{+}1)$.
\end{lemma}

\begin{proof}
\leanok
Matched cells retain their original values. The counting argument for the
parenthesis matching is preserved because the symmetric swap of unmatched
cells leaves the matching structure invariant.
\end{proof}

\begin{lemma}[BK content swap]
\label{lem.sf.bk-content-swap-lr}
\lean{AlgebraicCombinatorics.benderKnuth_content_swap}
\leanhelper
\leanok
Let $T' = \beta_k(T)$. Then
\[
  \operatorname{cont}(T')_k = \operatorname{cont}(T)_{k+1},
  \quad
  \operatorname{cont}(T')_{k+1} = \operatorname{cont}(T)_k,
  \quad
  \operatorname{cont}(T')_i = \operatorname{cont}(T)_i
  \text{ for } i \notin \{k, k{+}1\}.
\]
\end{lemma}

\begin{proof}
\leanok
For $i \notin \{k, k{+}1\}$, cells with value $i$ are unchanged by BK.
For the swap of $k$ and $k{+}1$ counts, decompose each set into forced cells,
matched free cells, and unmatched free cells. The forced bijection gives
equal forced counts; matched counts are equal by the matching definition;
and unmatched cells swap their values, contributing to the other count.
\end{proof}

\begin{lemma}[BK and the monomial of a tableau]
\label{lem.sf.bk-monomial}
\lean{AlgebraicCombinatorics.benderKnuth_monomial}
\leanhelper
\leanok
For the monomial associated to a tableau,
$x^{\operatorname{cont}(\beta_k(T))} = \operatorname{rename}_{(k\;k{+}1)}(x^{\operatorname{cont}(T)})$.
That is, the BK involution acts on monomials by the transposition $(k\;k{+}1)$.
\end{lemma}

\begin{proof}
\leanok
This follows directly from the content swap: renaming by $(k\;k{+}1)$
exchanges the exponents of $x_k$ and $x_{k+1}$.
\end{proof}

\subsection{Symmetry of skew Schur polynomials}

\begin{lemma}[Skew Schur polynomial is symmetric]
\label{thm.sf.skew-schur-symm-lr}
\lean{AlgebraicCombinatorics.skewSchurPoly_isSymmetric}
\leanhelper
\leanok
For $N$-partitions $\lambda$ and $\mu$, the skew Schur polynomial
$s_{\lambda/\mu} \in \mathcal{P}$ is symmetric (invariant under all
permutations of the variables).
\end{lemma}

\begin{proof}
\leanok
It suffices to show invariance under adjacent transpositions $(k\;k{+}1)$,
since these generate $S_N$.  For each $k$, the BK involution $\beta_k$ gives
a bijection on semistandard tableaux with
$x^{\operatorname{cont}(\beta_k(T))} = \operatorname{rename}_{(k\;k{+}1)}(x^{\operatorname{cont}(T)})$.
Reindexing the sum defining $s_{\lambda/\mu}$ by this bijection shows invariance.
\end{proof}

\begin{lemma}[Alternant times skew Schur expansion]
\label{lem.sf.alternant-mul-skew-schur}
\lean{AlgebraicCombinatorics.alternant_mul_skewSchurPoly_eq_sum}
\leanhelper
\leanok
For $N$-partitions $\lambda$, $\mu$ and any $N$-tuple $\alpha$:
\[
  a_\alpha \cdot s_{\lambda/\mu}
  = \sum_{T \in \mathrm{SSYT}(\lambda/\mu)} a_{\alpha + \operatorname{cont}(T)}.
\]
\end{lemma}

\begin{proof}
\leanok
Expand the alternant as $a_\alpha = \sum_\sigma (-1)^\sigma x^{\alpha \circ \sigma}$
and multiply by $s_{\lambda/\mu} = \sum_T x^{\operatorname{cont}(T)}$. Using
the symmetry of $s_{\lambda/\mu}$ and regrouping terms gives the result.
\end{proof}

\subsection{Stembridge's involution infrastructure}

\begin{definition}[Violator columns and misstep indices]
\label{def.sf.stemb-violator-misstep}
\lean{AlgebraicCombinatorics.violatorColumns, AlgebraicCombinatorics.isMisstep, AlgebraicCombinatorics.misstepSet}
\leanhelper
\leanok
Let $T$ be a semistandard tableau and $\nu$ an $N$-tuple. A positive integer
$j$ is a \emph{violator column} if $\nu + \operatorname{cont}(\operatorname{col}_{\geq j} T)$
is not an $N$-partition. An index $k$ is a \emph{misstep} for an $N$-tuple
$\alpha$ if $k + 1 < N$ and $\alpha_k < \alpha_{k+1}$ (i.e., the partition
condition fails at~$k$).
\end{definition}

\begin{lemma}[Non-Yamanouchi implies violator columns exist]
\label{lem.sf.stemb-violator-nonempty}
\lean{AlgebraicCombinatorics.violatorColumns_nonempty_of_not_yamanouchi}
\leanhelper
\leanok
If $T$ is semistandard and not $\nu$-Yamanouchi, then the set of violator
columns is nonempty.
\end{lemma}

\begin{proof}
\leanok
By definition, $T$ is not $\nu$-Yamanouchi iff there exists $j > 0$ such that
$\nu + \operatorname{cont}(\operatorname{col}_{\geq j} T)$ is not a partition.
\end{proof}

\begin{lemma}[Non-partition implies misstep exists]
\label{lem.sf.stemb-misstep-nonempty}
\lean{AlgebraicCombinatorics.not_isNPartition_iff_exists_misstep, AlgebraicCombinatorics.misstepSet_nonempty_of_not_partition}
\leanhelper
\leanok
An $N$-tuple $\alpha$ is not an $N$-partition if and only if it has a misstep.
In particular, the misstep set is nonempty when $\alpha$ is not a partition.
\end{lemma}

\begin{proof}
\leanok
If $\alpha$ is not weakly decreasing, there exist $i < j$ with
$\alpha_i < \alpha_j$. By well-founded induction on $j - i$, we find
consecutive indices with this property.
\end{proof}

\begin{definition}[Prefix-restricted Bender--Knuth]
\label{def.sf.bk-prefix}
\lean{AlgebraicCombinatorics.benderKnuthPrefixMatching}
\leanhelper
\leanok
The \emph{prefix-restricted BK involution} $\beta_k^{<j}$ applies the
parenthesis-matching swap of $k$ and $k{+}1$ only to cells in columns
$< j$. Cells in columns $\geq j$ are left unchanged:
\[
  \beta_k^{<j}(T)(c) = \begin{cases}
    k+1 & \text{if the column of } c \text{ is } < j \text{ and } c \text{ is unmatched free } k \text{ in prefix}, \\
    k   & \text{if the column of } c \text{ is } < j \text{ and } c \text{ is unmatched free } k{+}1 \text{ in prefix}, \\
    T(c) & \text{otherwise.}
  \end{cases}
\]
\end{definition}

\begin{lemma}[Prefix BK preserves suffix content]
\label{lem.sf.bk-prefix-content-colgeq}
\lean{AlgebraicCombinatorics.benderKnuthPrefixMatching_contentColGeq}
\leanhelper
\leanok
The prefix-restricted BK involution leaves columns $\geq j$ unchanged,
so
$\operatorname{cont}(\operatorname{col}_{\geq j} \beta_k^{<j}(T))
= \operatorname{cont}(\operatorname{col}_{\geq j} T)$.
\end{lemma}

\begin{proof}
\leanok
Cells in columns $\geq j$ are not modified by $\beta_k^{<j}$.
\end{proof}

\begin{lemma}[Prefix BK preserves semistandardness (Stembridge context)]
\label{lem.sf.bk-prefix-semistandard}
\lean{AlgebraicCombinatorics.benderKnuthPrefixMatching_semistandard_stembridge}
\leanhelper
\leanok
Under the Stembridge context (where $\beta = \nu + \operatorname{cont}(\operatorname{col}_{\geq j+1} T)$
is a partition and $k$ is a misstep for $\nu + \operatorname{cont}(\operatorname{col}_{\geq j} T)$),
the prefix-restricted BK involution $\beta_k^{<j}(T)$ preserves semistandardness.
\end{lemma}

\begin{proof}
\leanok
Row-weak and column-strict ordering are verified by case analysis,
similar to the full BK involution but restricted to the prefix.
The Stembridge context hypotheses ensure that boundary cells (at column $j$)
interact correctly with the modified prefix.
\end{proof}

\begin{lemma}[Prefix BK involutivity in Stembridge context]
\label{lem.sf.bk-prefix-involutive-stembridge}
\lean{AlgebraicCombinatorics.benderKnuthPrefixMatching_involutive_stembridge}
\leanhelper
\leanok
Under the Stembridge context, the prefix-restricted BK involution
is an involution: $\beta_k^{<j}(\beta_k^{<j}(T)) = T$. Moreover,
the Stembridge context is preserved: $T' = \beta_k^{<j}(T)$ again
satisfies the partition and misstep conditions needed for the next application.
\end{lemma}

\begin{proof}
\leanok
The suffix content is preserved (columns $\geq j$ are unchanged), so the
violator column and misstep index remain the same. Applying $\beta_k^{<j}$
twice returns every prefix cell to its original value by the matching symmetry.
\end{proof}

\begin{lemma}[Column $j$ has no entry $k$ at the max violator]
\label{lem.sf.column-j-no-k}
\lean{AlgebraicCombinatorics.column_j_no_k_at_max_violator}
\leanhelper
\leanok
Under the Stembridge context (where $j$ is the max violator column and $k$
is the misstep index), column $j$ contains no entry equal to $k$. In particular,
$\operatorname{cont}(\operatorname{col}_{\geq j} T)_k
= \operatorname{cont}(\operatorname{col}_{\geq j+1} T)_k$.
\end{lemma}

\begin{proof}
\leanok
Since $\beta = \nu + \operatorname{cont}(\operatorname{col}_{\geq j+1} T)$ is a partition
and $\alpha = \nu + \operatorname{cont}(\operatorname{col}_{\geq j} T)$ has a misstep at $k$,
the difference $\alpha_k - \beta_k$ must be zero. Otherwise the misstep at $k$
in $\alpha$ would contradict the partition property of $\beta$ combined with the
constraint $\alpha_{k+1} > \alpha_k$.
\end{proof}

\subsection{Stembridge's involution and sign-reversing property}

\begin{definition}[Stembridge's involution]
\label{def.sf.stemb-involution}
\lean{AlgebraicCombinatorics.stembridgeInvolution}
\leanhelper
\leanok
For a semistandard tableau $T$ that is not $\nu$-Yamanouchi, define
$T^* = \operatorname{stemb}_\nu(T)$ as follows:
\begin{enumerate}
  \item Let $j$ be the \emph{largest} violator column (where $\nu + \operatorname{cont}(\operatorname{col}_{\geq j} T)$
        is not a partition).
  \item Let $k$ be the \emph{smallest} misstep index of $\nu + \operatorname{cont}(\operatorname{col}_{\geq j} T)$.
  \item Apply the prefix-restricted BK involution: $T^* = \beta_k^{<j}(T)$.
\end{enumerate}
\end{definition}

\begin{lemma}[Stembridge involution preserves semistandardness]
\label{lem.sf.stemb-semistandard}
\lean{AlgebraicCombinatorics.stembridgeInvolution_semistandard}
\leanhelper
\leanok
If $T$ is semistandard and not $\nu$-Yamanouchi, then
$\operatorname{stemb}_\nu(T)$ is semistandard.
\end{lemma}

\begin{proof}
\leanok
Follows from the involutivity theorem, which establishes semistandardness
of $T'$ as a prerequisite.
\end{proof}

\begin{lemma}[Stembridge involution preserves non-Yamanouchi]
\label{lem.sf.stemb-not-yamanouchi}
\lean{AlgebraicCombinatorics.stembridgeInvolution_not_yamanouchi}
\leanhelper
\leanok
If $T$ is not $\nu$-Yamanouchi, then
$\operatorname{stemb}_\nu(T)$ is also not $\nu$-Yamanouchi.
\end{lemma}

\begin{proof}
\leanok
The violator column $j$ remains a violator for $T'$ because $\beta_k^{<j}$
leaves columns $\geq j$ unchanged, so
$\nu + \operatorname{cont}(\operatorname{col}_{\geq j} T')
= \nu + \operatorname{cont}(\operatorname{col}_{\geq j} T)$,
which is not a partition.
\end{proof}

\begin{lemma}[Max violator column is preserved]
\label{lem.sf.stemb-max-violator-preserved}
\lean{AlgebraicCombinatorics.max_violator_succ_isNPartition}
\leanhelper
\leanok
Let $j$ be the maximum violator column for $T$. Then
$\nu + \operatorname{cont}(\operatorname{col}_{\geq j+1} T)$ is an $N$-partition.
(This is because $j{+}1$ is \emph{not} a violator column, being larger than the maximum.)
If $j{+}1$ exceeds all column indices, the result reduces to $\nu$ being a partition.
\end{lemma}

\begin{proof}
\leanok
If $j{+}1 \leq \max_i \lambda_i + 1$, then $j{+}1$ is in the range of potential
violator columns but is not a violator (since $j$ is the max), so
$\nu + \operatorname{cont}(\operatorname{col}_{\geq j+1} T)$ is a partition
by definition. Otherwise, $\operatorname{cont}(\operatorname{col}_{\geq j+1} T) = \mathbf{0}$,
and the result is $\nu$, which is a partition by hypothesis.
\end{proof}

\begin{theorem}[Stembridge involution is an involution]
\label{lem.sf.stemb-involutive}
\lean{AlgebraicCombinatorics.stembridgeInvolution_involutive}
\leanhelper
\leanok
For $N$-partitions $\lambda, \mu, \nu$ and a semistandard tableau $T$ that
is not $\nu$-Yamanouchi: let $T' = \operatorname{stemb}_\nu(T)$. Then
$T'$ is semistandard, $T'$ is not $\nu$-Yamanouchi, and
$\operatorname{stemb}_\nu(T') = T$.
\end{theorem}

\begin{proof}
\leanok
The key insight is that the parameters $(j, k)$ are \emph{the same} for $T$
and $T'$: the max violator column $j'$ for $T'$ equals $j$ (since $\beta_k^{<j}$
leaves columns $\geq j$ unchanged, and columns $> j$ are not violators for either
tableau). The misstep set is unchanged because
$\operatorname{cont}(\operatorname{col}_{\geq j} T')
= \operatorname{cont}(\operatorname{col}_{\geq j} T)$.
Therefore $\operatorname{stemb}_\nu(T') = \beta_k^{<j}(\beta_k^{<j}(T)) = T$.
\end{proof}

\begin{lemma}[Stembridge involution content transposition]
\label{lem.sf.stemb-content-transposition}
\lean{AlgebraicCombinatorics.stembridgeInvolution_content_transposition}
\leanhelper
\leanok
Let $T' = \operatorname{stemb}_\nu(T)$. There exist indices $k \neq k'$
(namely $k$ and $k{+}1$) such that
\[
  \nu + \operatorname{cont}(T') + \rho
  = (\nu + \operatorname{cont}(T) + \rho) \circ (k\;k').
\]
That is, the full exponent tuple is related by a transposition.
\end{lemma}

\begin{proof}
\leanok
The proof uses the decomposition of content into prefix and suffix parts.
The suffix is unchanged ($\operatorname{col}_{\geq j}$ is preserved).
The prefix content swaps: the number of $k$'s in the prefix of $T'$ equals
the number of $(k{+}1)$'s in the prefix of $T$, and vice versa. Combined with
the misstep condition $\alpha_{k+1} = \alpha_k + 1$ and the rho shift
$\rho_k - \rho_{k+1} = 1$, the full expressions at positions $k$ and $k{+}1$
are exchanged.
\end{proof}

\begin{lemma}[Prefix BK content swap]
\label{lem.sf.bk-prefix-content-swap}
\lean{AlgebraicCombinatorics.benderKnuthPrefixMatching_content_swap}
\leanhelper
\leanok
For the prefix-restricted BK involution $T' = \beta_k^{<j}(T)$:
\[
  \operatorname{cont}(T')_k + \operatorname{cont}(\operatorname{col}_{\geq j} T)_{k+1}
  = \operatorname{cont}(T)_{k+1} + \operatorname{cont}(\operatorname{col}_{\geq j} T)_k.
\]
This identity captures how the BK swap in the prefix interacts with the
unchanged suffix.
\end{lemma}

\begin{proof}
\leanok
Decompose $\operatorname{cont}(T)_i = (\text{prefix } i \text{ count}) + \operatorname{cont}(\operatorname{col}_{\geq j} T)_i$.
The suffix is unchanged, so it suffices to show: $(\text{prefix } k \text{ count in } T') = (\text{prefix } (k{+}1) \text{ count in } T)$.
The prefix cells where $T'(c) = k$ are exactly the prefix cells that were $k$ in $T$
and stayed as $k$ (forced $k$'s and matched free $k$'s), plus the unmatched free
$(k{+}1)$'s (which became $k$).
The prefix cells where $T(c) = k{+}1$ decompose into cells that stay as $(k{+}1)$
(forced and matched) and unmatched $(k{+}1)$'s (which become $k$).
The unchanged cells contribute equally to both sides, and the
unmatched $(k{+}1)$'s appear on both sides as well, yielding the identity.
\end{proof}

\begin{lemma}[Stembridge involution is sign-reversing]
\label{lem.sf.stemb-sign-reversing}
\lean{AlgebraicCombinatorics.stembridgeInvolution_sign_reversing}
\leanhelper
\leanok
For a non-$\nu$-Yamanouchi semistandard tableau $T$:
\[
  a_{\nu + \operatorname{cont}(\operatorname{stemb}_\nu(T)) + \rho}
  = -\,a_{\nu + \operatorname{cont}(T) + \rho}.
\]
\end{lemma}

\begin{proof}
\leanok
By Lemma~\ref{lem.sf.stemb-content-transposition}, the exponent tuple is
related by a transposition. By Lemma~\ref{lem.sf.alternant-0-lr}~\textbf{(b)},
swapping two entries of the exponent negates the alternant.
\end{proof}

\begin{lemma}[Fixed points have zero alternant]
\label{lem.sf.stemb-fixed-point-zero}
\lean{AlgebraicCombinatorics.stembridgeInvolution_fixed_point_implies_alternant_zero}
\leanhelper
\leanok
If $\operatorname{stemb}_\nu(T) = T$ (a fixed point), then
$a_{\nu + \operatorname{cont}(T) + \rho} = 0$.
\end{lemma}

\begin{proof}
\leanok
If $T' = T$, then $\operatorname{cont}(T') = \operatorname{cont}(T)$, so the
content transposition gives $f = f \circ (k\;k')$ with $k \neq k'$, where
$f = \nu + \operatorname{cont}(T) + \rho$. This forces $f(k) = f(k')$, so
$f$ has repeated entries. By Lemma~\ref{lem.sf.alternant-0-lr}~\textbf{(a)},
$a_f = 0$.
\end{proof}

\begin{lemma}[Non-Yamanouchi alternants cancel]
\label{lem.sf.non-yamanouchi-cancel}
\lean{AlgebraicCombinatorics.alternant_sum_non_yamanouchi_eq_zero}
\leanhelper
\leanok
For $N$-partitions $\lambda, \mu, \nu$:
\[
  \sum_{\substack{T \in \mathrm{SSYT}(\lambda/\mu) \\ T \text{ not } \nu\text{-Yamanouchi}}}
  a_{\nu + \operatorname{cont}(T) + \rho} = 0.
\]
\end{lemma}

\begin{proof}
\leanok
The Stembridge involution partitions the non-Yamanouchi tableaux into orbits
of size $1$ or $2$. For each pair $\{T, T^*\}$ with $T^* \neq T$, the sign-reversing
property gives $a_{\nu + \operatorname{cont}(T) + \rho} + a_{\nu + \operatorname{cont}(T^*) + \rho} = 0$.
At fixed points (where $T^* = T$), the alternant is zero by
Lemma~\ref{lem.sf.stemb-fixed-point-zero}. Hence, all terms cancel and the sum is zero.
\end{proof}

\begin{lemma}[Alternant sum splits into Yamanouchi and non-Yamanouchi]
\label{lem.sf.alternant-sum-split}
\lean{AlgebraicCombinatorics.alternant_sum_split}
\leanhelper
\leanok
The sum of alternants over all semistandard tableaux equals the sum over
Yamanouchi tableaux plus the sum over non-Yamanouchi tableaux:
\[
  \sum_{T \in \mathrm{SSYT}} a_{\nu + \operatorname{cont}(T) + \rho}
  = \sum_{\substack{T \in \mathrm{SSYT} \\ \nu\text{-Yam}}} a_{\nu + \operatorname{cont}(T) + \rho}
  + \sum_{\substack{T \in \mathrm{SSYT} \\ \text{not } \nu\text{-Yam}}} a_{\nu + \operatorname{cont}(T) + \rho}.
\]
\end{lemma}

\begin{proof}
\leanok
The semistandard tableaux partition into Yamanouchi and non-Yamanouchi subsets.
The sum over a disjoint union equals the sum of the parts.
\end{proof}

\subsection{Stembridge's Lemma}

\begin{lemma}[Stembridge's Lemma]
\label{lem.sf.stemb-lem}
\lean{AlgebraicCombinatorics.stembridgeLemma}
\leantarget
\leanok
Let $\lambda,\mu,\nu$ be three
$N$-partitions. Then,%
\[
a_{\nu+\rho}\cdot s_{\lambda/\mu}=\sum_{\substack{T\text{ is a }%
\nu\text{-Yamanouchi}\\\text{semistandard tableau}\\\text{of shape }%
\lambda/\mu}}a_{\nu+\operatorname*{cont}T+\rho}.
\]
\end{lemma}

\begin{proof}
\leanok
The proof proceeds in several steps.

\textbf{Step 1.} By Lemma~\ref{lem.sf.alternant-mul-skew-schur}
(which uses the symmetry of $s_{\lambda/\mu}$ from
Lemma~\ref{thm.sf.skew-schur-symm-lr}):
\[
a_{\nu+\rho} \cdot s_{\lambda/\mu}
= \sum_{T \in \mathrm{SSYT}(\lambda/\mu)} a_{\nu + \operatorname{cont} T + \rho}.
\]

\textbf{Step 2.} By Lemma~\ref{lem.sf.alternant-sum-split}, split the sum into
Yamanouchi and non-Yamanouchi tableaux.

\textbf{Step 3.} By Lemma~\ref{lem.sf.non-yamanouchi-cancel}, the non-Yamanouchi
sum is zero (using the Stembridge involution from
Definition~\ref{def.sf.stemb-involution}, its sign-reversing property from
Lemma~\ref{lem.sf.stemb-sign-reversing}, and its fixed-point behavior from
Lemma~\ref{lem.sf.stemb-fixed-point-zero}).
\end{proof}

\subsection{Consequences of Stembridge's Lemma}

\begin{lemma}
\label{lem.sf.tab-greater-i}
\lean{AlgebraicCombinatorics.tableau_entry_ge_row}
\leantarget
\leanok
Let $\lambda$ be any $N$-partition. Let $T$ be a
semistandard tableau of shape $\lambda$. Then, $T\left(  i,j\right)  \geq i$
for each $\left(  i,j\right)  \in Y\left(  \lambda\right)  $.
\end{lemma}

\begin{proof}
\leanok
This lemma is an easy consequence
of the fact that the entries of a semistandard tableau increase strictly down
each column. More precisely, the cells
$(1, j), (2, j), \ldots, (i, j)$ all belong to $Y(\lambda)$ (since $\lambda$
is weakly decreasing), and the entries $T(1, j) < T(2, j) < \cdots < T(i, j)$
are strictly increasing. Since $T(1, j) \geq 1$, we get $T(i, j) \geq i$ by
induction.
\end{proof}

\begin{lemma}[Minimalistic tableau is semistandard]
\label{lem.sf.minimalistic-semistandard}
\lean{AlgebraicCombinatorics.minimalisticTableau_semistandard}
\leanhelper
\leanok
The minimalistic tableau $T_0$ of shape $\lambda$ (where $T_0(i,j) = i$
for all cells) is semistandard: entries are constant (hence weakly increasing)
along rows, and strictly increasing down columns.
\end{lemma}

\begin{proof}
\leanok
Row-weak: $T_0(i, j_1) = i = T_0(i, j_2)$ for any $j_1 < j_2$.
Column-strict: if $i_1 < i_2$, then $T_0(i_1, j) = i_1 < i_2 = T_0(i_2, j)$.
\end{proof}

\begin{lemma}[Minimalistic tableau is Yamanouchi]
\label{lem.sf.minimalistic-yamanouchi}
\lean{AlgebraicCombinatorics.minimalisticTableau_yamanouchi}
\leanhelper
\leanok
Let $\lambda$ be an $N$-partition. The minimalistic tableau $T_0$ of shape
$\lambda$ (whose entry at each cell $(i,j)$ is $i$) is a $\mathbf{0}$-Yamanouchi semistandard tableau of shape
$\lambda/\mathbf{0}$.
\end{lemma}

\begin{proof}
\leanok
The minimalistic tableau is semistandard (entries weakly increase along rows
and strictly increase down columns). For the $\mathbf{0}$-Yamanouchi condition,
the content of $\operatorname{col}_{\geq j} T_0$ at row $i$ equals
$\max\{\lambda_i - (j-1), 0\}$, which is weakly decreasing since $\lambda$ is.
\end{proof}

\begin{lemma}[Content of minimalistic tableau]
\label{lem.sf.minimalistic-content}
\lean{AlgebraicCombinatorics.contentTableau_minimalisticTableau}
\leanhelper
\leanok
Let $\lambda$ be an $N$-partition. The content of the minimalistic tableau $T_0$
equals $\lambda$: $\operatorname{cont}(T_0) = \lambda$.
\end{lemma}

\begin{proof}
\leanok
The minimalistic tableau has all entries in row $i$ equal to $i$, so the number
of $i$'s in $T_0$ is $\lambda_i$ (the number of cells in row $i$).
\end{proof}

\begin{lemma}[Yamanouchi rightmost entry equals row index]
\label{lem.sf.yamanouchi-rightmost}
\lean{AlgebraicCombinatorics.yamanouchi_rightmost_entry}
\leanhelper
\leanok
Let $\lambda$ be an $N$-partition and $T$ a $\mathbf{0}$-Yamanouchi tableau
of shape $\lambda/\mathbf{0}$. For each row $i$ with $\lambda_i > 0$,
the rightmost cell $(i, \lambda_i)$ satisfies $T(i, \lambda_i) = i$.
\end{lemma}

\begin{proof}
\leanok
By strong induction on $i$. Suppose $T(i, \lambda_i) = m \neq i$.
By Lemma~\ref{lem.sf.tab-greater-i}, $m \geq i$, and since $m \neq i$, we have $m > i$.
From the $\mathbf{0}$-Yamanouchi condition,
$\operatorname{cont}(\operatorname{col}_{\geq \lambda_i} T)$ is a partition,
so $\operatorname{cont}(\operatorname{col}_{\geq \lambda_i} T)_i
\geq \operatorname{cont}(\operatorname{col}_{\geq \lambda_i} T)_m \geq 1$.
Hence row $i$ has some entry $i$ in columns $\geq \lambda_i$.
But $(i, \lambda_i)$ is the rightmost cell of row $i$, so
$T(i, \lambda_i) \leq i$ by weak row increase from the cell where $i$ appears.
For rows $i' < i$, the inductive hypothesis gives $T(i', \lambda_{i'}) = i'$;
the strict column increase then prevents any cell in the same column as
$(i, \lambda_i)$ from having value $i$ in an earlier row.
This contradicts $m > i$.
\end{proof}

\begin{lemma}[Uniqueness of $\mathbf{0}$-Yamanouchi tableau]
\label{lem.sf.minimalistic-unique}
\lean{AlgebraicCombinatorics.minimalisticTableau_unique}
\leanhelper
\leanok
Let $\lambda$ be an $N$-partition. If $T$ is a $\mathbf{0}$-Yamanouchi
semistandard tableau of shape $\lambda/\mathbf{0}$, then $T$ is the
minimalistic tableau $T_0$.
\end{lemma}

\begin{proof}
\leanok
By Lemma~\ref{lem.sf.yamanouchi-rightmost}, $T(i, \lambda_i) = i$ for each row.
By Lemma~\ref{lem.sf.tab-greater-i}, $T(i, j) \geq i$ for all $(i, j)$.
By weak row increase, $T(i, j) \leq T(i, \lambda_i) = i$ for all $j \leq \lambda_i$.
Combining: $T(i, j) = i$ for all $(i, j)$, so $T = T_0$.
\end{proof}

\begin{lemma}[Schur--alternant relation]
\label{lem.sf.schur-alternant}
\lean{AlgebraicCombinatorics.schurPoly_eq_alternant_div}
\leanhelper
\leanok
Let $\lambda$ be an $N$-partition. Then,
\[
a_{\lambda + \rho} = a_{\rho} \cdot s_{\lambda}.
\]
(This is Theorem~\ref{thm.sf.schur-symm}~\textbf{(b)}.)
\end{lemma}

\begin{proof}
\leanok
Apply Lemma~\ref{lem.sf.stemb-lem} with $\mu = \mathbf{0}$ and
$\nu = \mathbf{0}$:
\[
a_{\rho} \cdot s_{\lambda}
= \sum_{\substack{T \text{ is a } \mathbf{0}\text{-Yamanouchi} \\
    \text{semistandard tableau of shape } \lambda/\mathbf{0}}}
  a_{\operatorname{cont} T + \rho}.
\]
By Lemma~\ref{lem.sf.minimalistic-yamanouchi}, the minimalistic tableau $T_0$
is $\mathbf{0}$-Yamanouchi, and by Lemma~\ref{lem.sf.minimalistic-unique},
it is the \emph{only} such tableau. Therefore the sum has a single term.
By Lemma~\ref{lem.sf.minimalistic-content},
$\operatorname{cont}(T_0) = \lambda$, so we obtain
$a_{\rho} \cdot s_{\lambda} = a_{\lambda + \rho}$.
\end{proof}

\subsection{The Littlewood--Richardson rule}

\begin{theorem}[Zelevinsky's generalized Littlewood--Richardson rule, in Yamanouchi form]
\label{thm.sf.lr-zy}
\lean{AlgebraicCombinatorics.littlewoodRichardson}
\leantarget
\leanok
Let $\lambda,\mu,\nu$ be three $N$-partitions. Then,%
\begin{equation}
s_{\nu}\cdot s_{\lambda/\mu}=\sum_{\substack{T\text{ is a }\nu
\text{-Yamanouchi}\\\text{semistandard tableau}\\\text{of shape }\lambda/\mu
}}s_{\nu+\operatorname*{cont}T}. \label{eq.thm.sf.lr-zy.eq}%
\end{equation}
\end{theorem}

\begin{proof}
\leanok
\textbf{Step 1.} By Lemma~\ref{lem.sf.schur-alternant} (applied to $\nu$
instead of $\lambda$), $a_{\nu+\rho} = a_{\rho} \cdot s_{\nu}$.

\textbf{Step 2.} By Lemma~\ref{lem.sf.stemb-lem},
\[
a_{\nu+\rho} \cdot s_{\lambda/\mu}
= \sum_{\substack{T \text{ is a } \nu\text{-Yamanouchi} \\
    \text{semistandard tableau} \\ \text{of shape } \lambda/\mu}}
  a_{\nu + \operatorname{cont} T + \rho}.
\]

\textbf{Step 3.} For each $\nu$-Yamanouchi tableau $T$, the tuple
$\nu + \operatorname{cont} T$ is an $N$-partition (by the $j=1$ case of the
Yamanouchi condition), so Lemma~\ref{lem.sf.schur-alternant} yields
$a_{\nu + \operatorname{cont} T + \rho} = a_{\rho} \cdot s_{\nu + \operatorname{cont} T}$.

\textbf{Step 4.} Substituting Steps 1 and 3 into Step 2:
\[
a_{\rho} \cdot s_{\nu} \cdot s_{\lambda/\mu}
= a_{\rho} \cdot \sum_{\substack{T \text{ is a } \nu\text{-Yamanouchi} \\
    \text{semistandard tableau} \\ \text{of shape } \lambda/\mu}}
  s_{\nu + \operatorname{cont} T}.
\]

\textbf{Step 5.} Since $a_{\rho}$ is regular (Lemma~\ref{lem.sf.arho-reg}),
we can cancel it (Lemma~\ref{lem.cring.reg.cancel}) to obtain
\[
s_{\nu} \cdot s_{\lambda/\mu}
= \sum_{\substack{T \text{ is a } \nu\text{-Yamanouchi} \\
    \text{semistandard tableau} \\ \text{of shape } \lambda/\mu}}
  s_{\nu + \operatorname{cont} T}.
\]
\end{proof}
